% Part II -- Python Foundations
% Chapter 2: Values, Variables, and Calculations

\h{Values, Variables, and Calculations}

A useful program needs information to work with. That information may be a
number, a piece of text, or a yes-or-no value. In this chapter, you will learn
how Python represents these values, how to give them clear names, and how to
use them in calculations.

\begin{NoteBox}
\textbf{Prerequisites.} You should be able to create, save, and run a Python
file and use \c{print()}.
\end{NoteBox}

\begin{tcolorbox}[title=Learning outcomes,colframe=orange,colback=white,arc=2mm]
After completing this chapter, you should be able to:
\begin{itemize}
    \item recognize integers, floating-point numbers, strings, booleans, and \c{None};
    \item assign values to clearly named variables;
    \item perform and organize arithmetic calculations;
    \item read text with \c{input()} and convert it to a number; and
    \item produce readable output with an f-string.
\end{itemize}
\end{tcolorbox}

\hh{Values and types}

A \emph{value} is one piece of information. Every value has a \emph{type}. The
type tells Python what kind of information the value represents and which
operations make sense for it.

\begin{lstlisting}[
    language=python,
    style=block,
    caption={Five basic Python values},
    label={c02_basic_values},
    captionpos=b
]
print(4)
print(3.5)
print("Studio")
print(True)
print(None)
\end{lstlisting}

These values have different types:
\begin{itemize}
    \item \c{4} is an integer, written \c{int}.
    \item \c{3.5} is a floating-point number, written \c{float}.
    \item \c{"Studio"} is a string, written \c{str}.
    \item \c{True} is a boolean, written \c{bool}.
    \item \c{None} represents the absence of a value.
\end{itemize}

Use \c{type()} when you want Python to report a value's type.

\begin{lstlisting}[
    language=python,
    style=block,
    caption={Inspect three value types},
    label={c02_type_check},
    captionpos=b
]
print(type(4))
print(type(3.5))
print(type("Studio"))
\end{lstlisting}

\begin{NoteBox}
\textbf{Stop and predict.} What type will Python report for \c{"4"}? Notice
that quotation marks change a number into text.
\end{NoteBox}

\hh{Variables give values names}

A variable is a name connected to a value. The equals sign assigns the value
on the right to the name on the left.

\begin{lstlisting}[
    language=python,
    style=block,
    caption={Assign and display three variables},
    label={c02_first_variables},
    captionpos=b
]
room_name = "Studio"
width_m = 6.0
length_m = 8.0

print(room_name)
print(width_m)
print(length_m)
\end{lstlisting}

Read \c{width_m = 6.0} as ``assign the value 6.0 to \c{width_m}.'' The equals
sign performs assignment here; it is not asking a mathematical question.

Use names that explain the value:
\begin{itemize}
    \item begin with a letter or underscore;
    \item use letters, numbers, and underscores only;
    \item use lowercase words separated by underscores; and
    \item include units when they prevent confusion.
\end{itemize}

Names such as \c{wall_height_m} and \c{area_m2} are clearer than \c{x} and
\c{result}.

\hh{Variables can be reassigned}

A variable can receive a new value later.

\begin{lstlisting}[
    language=python,
    style=block,
    caption={Reassign a variable},
    label={c02_reassignment},
    captionpos=b
]
floor_count = 3
print(floor_count)

floor_count = 4
print(floor_count)
\end{lstlisting}

Expected output:

\begin{lstlisting}[
    language=text,
    style=block,
    caption={Output after reassignment},
    label={c02_reassignment_output},
    captionpos=b
]
3
4
\end{lstlisting}

Reassignment does not change the earlier output. It changes which value the
name refers to from that point forward.

\hh{Work with numbers}

Python uses familiar arithmetic operators.

\begin{center}
\begin{tabular}{lll}
\textbf{Operation} & \textbf{Operator} & \textbf{Example} \\
Addition & \c{+} & \c{8 + 2} \\
Subtraction & \c{-} & \c{8 - 2} \\
Multiplication & \c{*} & \c{8 * 2} \\
Division & \c{/} & \c{8 / 2} \\
Floor division & \c{//} & \c{9 // 2} \\
Exponentiation & \c{**} & \c{3 ** 2} \\
\end{tabular}
\end{center}

Parentheses make the intended order clear.

\begin{lstlisting}[
    language=python,
    style=block,
    caption={Calculate area and volume},
    label={c02_arithmetic},
    captionpos=b
]
width_m = 6.0
length_m = 8.0
height_m = 3.2

area_m2 = width_m * length_m
volume_m3 = area_m2 * height_m

print(area_m2)
print(volume_m3)
\end{lstlisting}

\begin{NoteBox}
\textbf{Stop and predict.} If \c{width_m} changes to \c{7.0}, which results
will change? Make a prediction before running the revised code.
\end{NoteBox}

\hh{Control calculation order}

Python follows standard arithmetic precedence. Multiplication and division
occur before addition and subtraction. Parentheses let you state the intended
order directly.

\begin{lstlisting}[
    language=python,
    style=block,
    caption={Parentheses change calculation order},
    label={c02_precedence},
    captionpos=b
]
result_a = 4 + 3 * 2
result_b = (4 + 3) * 2

print(result_a)
print(result_b)
\end{lstlisting}

Expected output:

\begin{lstlisting}[
    language=text,
    style=block,
    caption={Results with and without parentheses},
    label={c02_precedence_output},
    captionpos=b
]
10
14
\end{lstlisting}

Use parentheses when a reader might otherwise need to stop and check the
order.

\hh{Work with strings}

A string is text inside quotation marks. Single and double quotation marks
both create strings.

\begin{lstlisting}[
    language=python,
    style=block,
    caption={Create two strings},
    label={c02_strings},
    captionpos=b
]
room_name = "Lecture Room"
level_name = 'Second Floor'

print(room_name)
print(level_name)
\end{lstlisting}

Choose one quotation style and use it consistently. A string can contain
spaces, punctuation, and digits. The digits in \c{"Room 204"} are still text.

Strings can be joined with \c{+}.

\begin{lstlisting}[
    language=python,
    style=block,
    caption={Join two strings},
    label={c02_string_join},
    captionpos=b
]
room_name = "Studio"
message = "Room: " + room_name
print(message)
\end{lstlisting}

For messages containing numbers, an f-string is usually clearer than joining
many converted values.

\hh{Read input from the student}

The \c{input()} function displays a prompt, waits for the user to type, and
returns the response as a string.

\begin{lstlisting}[
    language=python,
    style=block,
    caption={Read and display text input},
    label={c02_text_input},
    captionpos=b
]
room_name = input("Room name: ")
print("You entered: " + room_name)
\end{lstlisting}

Even when a user types digits, \c{input()} returns text. Convert the response
before using it in arithmetic.

\begin{lstlisting}[
    language=python,
    style=block,
    caption={Convert input text to a floating-point number},
    label={c02_numeric_input},
    captionpos=b
]
width_text = input("Width in meters: ")
width_m = float(width_text)

print(width_m * 2)
\end{lstlisting}

You can combine reading and conversion when the shorter form remains clear.

\begin{lstlisting}[
    language=python,
    style=block,
    caption={Read two dimensions and calculate area},
    label={c02_direct_conversion},
    captionpos=b
]
width_m = float(input("Width in meters: "))
length_m = float(input("Length in meters: "))
area_m2 = width_m * length_m

print(area_m2)
\end{lstlisting}

If the user types text that cannot be converted, Python reports a
\c{ValueError}. Chapter 8 will show how to handle that error. For now, enter a
valid number when a prompt requests one.

\hh{Format output with f-strings}

An f-string places values inside a readable message. Begin the string with
\c{f} and place variable names inside braces.

\begin{lstlisting}[
    language=python,
    style=block,
    caption={Insert values into an f-string},
    label={c02_first_fstring},
    captionpos=b
]
room_name = "Studio"
area_m2 = 48.0

print(f"{room_name} area: {area_m2} m2")
\end{lstlisting}

Use a format specification to control decimal places.

\begin{lstlisting}[
    language=python,
    style=block,
    caption={Display a number with two decimal places},
    label={c02_decimal_format},
    captionpos=b
]
width_m = 5.75
length_m = 7.2
area_m2 = width_m * length_m

print(f"Area: {area_m2:.2f} m2")
\end{lstlisting}

The specification \c{:.2f} displays a floating-point value with two digits
after the decimal point. It changes the displayed form, not the stored value.

\hh{Use comments to explain intent}

A comment begins with \c{#}. Python ignores the comment when the program runs.
Use comments to explain an assumption or decision.

\begin{lstlisting}[
    language=python,
    style=block,
    caption={A comment records an important assumption},
    label={c02_comments},
    captionpos=b
]
# Dimensions are measured in meters.
width_m = 5.75
length_m = 7.2

area_m2 = width_m * length_m
print(f"Area: {area_m2:.2f} m2")
\end{lstlisting}

Avoid comments that merely repeat the code. A comment should add information
that the names and operations do not already communicate.

\hh{Common mistakes}

\hhh{Using a variable before assigning it}

\begin{lstlisting}[
    language=python,
    style=block,
    caption={This intentionally produces a NameError},
    label={c02_name_error},
    captionpos=b
]
print(area_m2)
\end{lstlisting}

Assign \c{area_m2} before trying to use it.

\hhh{Mixing text and numbers with plus}

\begin{lstlisting}[
    language=python,
    style=block,
    caption={This intentionally produces a TypeError},
    label={c02_type_error},
    captionpos=b
]
area_m2 = 48.0
print("Area: " + area_m2)
\end{lstlisting}

Use an f-string instead:

\begin{lstlisting}[
    language=python,
    style=block,
    caption={Use an f-string to combine text and a number},
    label={c02_type_error_fixed},
    captionpos=b
]
area_m2 = 48.0
print(f"Area: {area_m2} m2")
\end{lstlisting}

\hhh{Using an unclear name}

\begin{lstlisting}[
    language=python,
    style=block,
    caption={Valid code with unclear names},
    label={c02_unclear_names},
    captionpos=b
]
a = 6.0
b = 8.0
c = a * b
\end{lstlisting}

Prefer \c{width_m}, \c{length_m}, and \c{area_m2}. Clear names reduce the
amount of explanation a reader needs.

\hh{Practice ladder}

\hhh{Level 0 -- Read and predict}

Predict the output before running the code.

\begin{lstlisting}[
    language=python,
    style=block,
    caption={Predict a calculation result},
    label={c02_practice_level0},
    captionpos=b
]
width_m = 4.0
length_m = 6.0
area_m2 = width_m * length_m
print(area_m2)
\end{lstlisting}

\hhh{Level 1 -- Modify}

Run the code above. Change the width and length to two different values. Check
the result by hand, then run the code again.

\hhh{Level 2 -- Complete}

Replace \c{...} with the missing expression so the program calculates
perimeter. Complete the line before you run the program.

\begin{lstlisting}[
    language=python,
    style=block,
    caption={Complete one missing expression},
    label={c02_practice_level2},
    captionpos=b
]
width_m = 5.0
length_m = 9.0
perimeter_m = ...  # Replace ... with the perimeter calculation.

print(f"Perimeter: {perimeter_m:.1f} m")
\end{lstlisting}

\hhh{Level 3 -- Apply}

Write a program that asks for width, length, and height. Calculate area and
volume, then display both results with two decimal places.

\textbf{Hints}
\begin{itemize}
    \item \textbf{Hint 1}: Use \c{float(input(...))} three times.
    \item \textbf{Hint 2}: Calculate area before volume.
    \item \textbf{Hint 3}: Volume equals area multiplied by height.
\end{itemize}

\hhh{Level 4 -- Challenge}

Extend Level 3 by calculating the difference between the length and width.
Display a clear sentence describing the difference.

\hh{Practice solutions}

\textbf{Level 0:} The output is \c{24.0}.

\textbf{Level 2}

\begin{lstlisting}[
    language=python,
    style=block,
    caption={One solution to Level 2},
    label={c02_solution_level2},
    captionpos=b
]
width_m = 5.0
length_m = 9.0
perimeter_m = 2 * (width_m + length_m)

print(f"Perimeter: {perimeter_m:.1f} m")
\end{lstlisting}

\textbf{Level 3}

\begin{lstlisting}[
    language=python,
    style=block,
    caption={One solution to Level 3},
    label={c02_solution_level3},
    captionpos=b
]
width_m = float(input("Width in meters: "))
length_m = float(input("Length in meters: "))
height_m = float(input("Height in meters: "))

area_m2 = width_m * length_m
volume_m3 = area_m2 * height_m

print(f"Area: {area_m2:.2f} m2")
print(f"Volume: {volume_m3:.2f} m3")
\end{lstlisting}

\textbf{Level 4 addition}

\begin{lstlisting}[
    language=python,
    style=block,
    caption={One solution that extends Level 3},
    label={c02_solution_level4},
    captionpos=b
]
width_m = float(input("Width in meters: "))
length_m = float(input("Length in meters: "))
height_m = float(input("Height in meters: "))

area_m2 = width_m * length_m
volume_m3 = area_m2 * height_m
difference_m = length_m - width_m

print(f"Area: {area_m2:.2f} m2")
print(f"Volume: {volume_m3:.2f} m3")
print(f"Length minus width: {difference_m:.2f} m")
\end{lstlisting}

\hh{Chapter checkpoint}

\begin{enumerate}
    \item What is the difference between \c{4} and \c{"4"}?
    \item What does the assignment operator do?
    \item Why is \c{area_m2} clearer than \c{x}?
    \item Which operator performs multiplication?
    \item Why might you use parentheses in a calculation?
    \item What type does \c{input()} return?
    \item How do you convert \c{"3.5"} to a number?
    \item What does the \c{f} before an f-string mean for its braces?
    \item What does \c{:.2f} control?
    \item What symbol begins a comment?
\end{enumerate}

\hh{Checkpoint answers}

\begin{enumerate}
    \item \c{4} is an integer; \c{"4"} is a string.
    \item It connects a name to a value.
    \item The name describes the quantity and its unit.
    \item \c{*}.
    \item To make the intended order clear or change the normal order.
    \item A string.
    \item Use \c{float("3.5")}.
    \item Values or expressions inside braces are inserted into the text.
    \item The displayed number of decimal places.
    \item \c{#}.
\end{enumerate}

\begin{tcolorbox}[title=Chapter summary,colframe=orange,colback=white,arc=2mm]
Values have types. Variables give values meaningful names. Python can perform
calculations with numeric variables, read responses with \c{input()}, convert
text to numbers, and insert results into f-strings. Chapter 3 shows how to
organize several related values in collections.
\end{tcolorbox}


